\section*{Limitations}

\paragraph{No human validation this cycle.} Coding reliability rests on two LLM judges from
different model families. Agreement is $\kappa=\Kappa{}$ on the main chains but only
$\KappaEShift{}$ on the \textsc{e-shift} target items, which are the hardest cases by construction,
and \EShiftUnresRate{} of those are unresolved. We therefore report every headline conclusion under
each judge separately and bound the unresolved cases in both directions (\S\ref{sec:c1}). Human
adjudication would be the stronger check; we have archived a stratified 150-item sample for it, and it
is not used in any inference reported here.

\paragraph{Prior-fit is not identified in isolation.} This is the paper's main limitation and it is
ours, not the field's. The frames move the compactor's elicited prior strongly and in the predicted
direction for every family, but a blind check (\S\ref{sec:parity}) shows they also lower judged
importance and authority, in the direction that would depress retention of the expected-value arm. A
real prior-fit benefit could therefore be masked by an importance penalty. \textsc{E-Shift} bounds a
compound context manipulation; it is an informative failed test of the mechanism, not a refutation of
it. The frame-free dose--response carries no such confound but is ecological and imprecise
(\DoseIQR{}~pp over the IQR, 95\% CI [\DoseIQRlo{}, \DoseIQRhi{}]). A manipulation that moves
expectation without moving perceived importance is the obvious next instrument, and we do not have
one.

\paragraph{The effect size we could test is larger than the one we wanted.} Between-family
heterogeneity in the frame effect is large and is not reducible by adding seeds, so the operative
smallest effect of interest had to be raised from the \(8\)~pp we preregistered to
\(\SESOI{}\)~pp. We do not claim equivalence at \(8\)~pp, and reaching \(0.80\) power there would
need roughly \(60\) eligible families.

\paragraph{The operator comparisons are exploratory.} Only \textsc{E-Shift} and the
surprise\,$\times$\,operator contrast were preregistered as confirmatory. The contrasts in
\S\ref{sec:operators} and \S\ref{sec:c7} were not, and are reported without multiplicity
correction. \textsc{Src} differs from \textsc{Roll} in more than iteration --- repeated exposure,
input length and composition, and recoverability --- so it identifies reconstruction-versus-recursion.
\textsc{Anchor} is a \emph{pipeline-level} comparison: it holds the storage cap fixed but not the
generation process (\RetriesOfive{} budget regenerations and \TrimsOfive{} trims per trajectory
against \RetriesOone{} and \TrimsOone{}), and since more trimming is associated with \emph{higher}
retention we cannot treat that difference as a handicap. Its gain is not attributable to the
anti-drift wording alone.

\paragraph{Semantic prior and representational compressibility are not fully separable.} Values a
model finds expected may also be easier to phrase compactly. \textsc{E-Shift} holds the surface form
exactly fixed, and the prior-blind controls remove model-driven selection, which constrains this
substantially; but both may be channels through which model expectations act, and we do not claim to
have separated them. Relatedly, our prior-blind controls trade off against each other: the one matched
on retention carries more text than the rewriting arm, and the one matched on the word budget retains
almost nothing. Neither alone is a perfect counterfactual, which is why both are reported.

\paragraph{The operative prior may move with $k$.} We measure priors at seeding and re-elicit them in
masked contexts at $k{=}1$ and $k{=}10$ on a diagnostic subset, reporting rank stability. A fully
time-varying dose is out of scope.

\paragraph{Generalisation.} The confirmatory result is for one strong compactor at one compression
regime. The second strong family and the small local model are \emph{portability diagnostics} on the
graded dose--response, not repetitions of \textsc{E-Shift}, and they score each model against the
primary compactor's measured prior rather than its own. Any null is reported as scoped to the tested
model, prompt, item domain and compression regime \citep{lakens2017equivalence}, never as a general
statement about compactors.

\paragraph{Simulated users, and a descriptive preference result.} The eight preference profiles are
controlled test cases, not a sample of people; generalisation is over preference families. The
non-default survival advantage is a robust descriptive association --- it holds within families as well
as across them --- but default-matching and non-default levels differ in wording and difficulty too, so
we do not attach it causally to the prior mechanism.

\section*{Ethical considerations}

This work studies how agent memory systems lose information, using entirely synthetic stimuli set in
a fictional workplace. No human subjects were involved and no personal data was collected or
processed. The falsehood-free design matters here: every stimulus value is true within its world, so
nothing we measure depends on the model being told something false, and nothing in this paper
describes a way to make a memory system behave badly on purpose. The practical upshot --- that
rewriting memory erodes exactly the user-specific content that personalisation depends on, and that
different users' stored profiles converge --- is a reason to design memory systems that keep what
users actually told them. All model calls were made through a budgeted proxy; total inference cost is
reported in Appendix~\ref{app:cost}.
